%表紙
\newpage
\chapter{講義規則}
\begin{enumerate}[label=第\arabic*条　]
  \item この章は、音楽ゲーム学園学園憲章（以下「憲章」という。）の趣旨に基づき、講義の実施に関し必要な事項を定めるものとする。
  \item 講師は、所定の申請を経て、講義を開講することができる。
    \begin{enumerate}[label=\arabic*　,start=2,leftmargin=*]
      \item 前項の申請の具体的な手続については、別に定める。
    \end{enumerate}
  \item 学生は、学園の定めるところにより、講義を受講することができる。
    \begin{enumerate}[label=\arabic*　,start=2,leftmargin=*]
      \item 前項の規定にかかわらず、講師は、担当する講義について、受講に関する条件を定めることができる。
    \end{enumerate}
  \item 講義資料の著作権その他の知的財産権は、当該資料を作成した講師に帰属する。
    \begin{enumerate}[label=\arabic*　,start=2,leftmargin=*]
      \item 前項の規定にかかわらず、学園は、第４章第４条第１号の同意に基づき、学内活動の目的に必要な範囲において、講義資料を無償で使用し、及び配信することができる。
    \end{enumerate}
  \item 講師は、あらかじめ予定した講義の回数を正規活動期間内に実施することが難しい場合、第２章に定める補講期間内に、その不足する回数を実施することができる。
  \item 学園長又は教務主事は、講義について次の各号のいずれかに該当する事実があると認めるときに限り、当該講師に対し、その中止又は変更を求めることができる。
      \begin{enumerate}[label=(\arabic*)　]
        \item 明らかに虚偽の情報の伝達その他講義の真正性を著しく損なう事実があるとき
        \item 受講者その他の学生に対するハラスメントその他著しく不適切な言動があるとき
        \item 前二号に掲げるもののほか、学生の権利利益を著しく損なうおそれがあると認められる事実があるとき
      \end{enumerate}
    \begin{enumerate}[label=\arabic*　,start=2,leftmargin=*]
      \item 前項の求めは、憲章の定める探究の自由を不当に制限することのないよう、必要な最小限度にとどめなければならない。
    \end{enumerate}
  \item この章の実施に関し必要な事項は、学園長が別に定める。
\end{enumerate}